\documentclass[preview]{standalone}

\usepackage{amsmath}
\usepackage{amssymb}
\usepackage{stellar}
\usepackage{definitions}
\usepackage{bettelini}

\begin{document}

\id{probability-gaussian-vectors}
\genpage

\section{Real Gaussian laws}

\begin{snippetdefinition}{gaussian-distribution-definition}{Gaussian distribution}
    Let \(m\in\realnumbers\) and let \(\sigma>0\). A real-valued
    \randomvariable \(X\) has the \emph{Gaussian distribution}, or
    \emph{normal distribution}, with mean \(m\) and variance \(\sigma^2\),
    written
    \[
        X\sim\mathcal N(m,\sigma^2),
    \]
    if \(X\) is an \absolutelycontinuousrandomvariable with density
    \[
        f_X(x)
        =
        \frac{1}{\sqrt{2\pi\sigma^2}}
        e^{-\frac{(x-m)^2}{2\sigma^2}}.
    \]
\end{snippetdefinition}

\begin{snippetproposition}{gaussian-distribution-affine-standardization}{Affine standardization of a Gaussian random variable}
    Let \(m\in\realnumbers\), let \(\sigma>0\), and let
    \(Z\sim\mathcal N(0,1)\). Then
    \[
        m+\sigma Z\sim\mathcal N(m,\sigma^2).
    \]
    Conversely, if \(X\sim\mathcal N(m,\sigma^2)\), then
    \[
        \frac{X-m}{\sigma}\sim\mathcal N(0,1).
    \]
\end{snippetproposition}

\begin{snippetproof}{gaussian-distribution-affine-standardization-proof}{gaussian-distribution-affine-standardization}{Affine standardization of a Gaussian random variable}
    The first statement follows from the one-dimensional change of variables
    \(x=m+\sigma z\):
    \[
        f_{m+\sigma Z}(x)
        =
        \frac{1}{\sigma}
        \frac{1}{\sqrt{2\pi}}
        e^{-\frac12\left(\frac{x-m}{\sigma}\right)^2}
        =
        \frac{1}{\sqrt{2\pi\sigma^2}}
        e^{-\frac{(x-m)^2}{2\sigma^2}}.
    \]
    The converse is the same computation applied to
    \(z=(x-m)/\sigma\).
\end{snippetproof}

\begin{snippetproposition}{standard-gaussian-moments}{Moments of the standard Gaussian distribution}
    Let \(Z\sim\mathcal N(0,1)\). Then, for every \(k\in\naturalnumbers\),
    \[
        \expectation[Z^{2k+1}]=0
    \]
    and
    \[
        \expectation[Z^{2k}]
        =
        (2k-1)(2k-3)\cdots 3\cdot1,
    \]
    where the empty product for \(k=0\) is \(1\).
\end{snippetproposition}

\begin{snippetproof}{standard-gaussian-moments-proof}{standard-gaussian-moments}{Moments of the standard Gaussian distribution}
    The odd moments vanish because \(z^{2k+1}e^{-z^2/2}\) is an odd
    function. For the even moments, define
    \[
        m_{2k}
        =
        \expectation[Z^{2k}]
        =
        \frac{1}{\sqrt{2\pi}}
        \int_{\realnumbers}z^{2k}e^{-\frac{z^2}{2}}\,dz.
    \]
    For \(k\geq1\), integration by parts gives
    \[
        \int_{\realnumbers}z^{2k}e^{-\frac{z^2}{2}}\,dz
        =
        (2k-1)
        \int_{\realnumbers}z^{2k-2}e^{-\frac{z^2}{2}}\,dz.
    \]
    Hence
    \[
        m_{2k}=(2k-1)m_{2k-2}.
    \]
    Since \(m_0=1\), iterating the recurrence gives
    \[
        m_{2k}=(2k-1)(2k-3)\cdots3\cdot1.
    \]
\end{snippetproof}

\section{Mean vectors and covariance matrices}

\begin{snippetdefinition}{expected-vector-definition}{Expected vector}
    Let \((\Omega,\mathcal F,\mathbb P)\) be a \probabilityspace, and let
    \[
        X=(X_1,\ldots,X_n)
    \]
    be a \randomvector in \(\realnumbers^n\). If each component \(X_i\) has
    a finite \expectedvalue, the \emph{expected vector} of \(X\) is
    \[
        \expectation[X]
        =
        \begin{pmatrix}
            \expectation[X_1]\\
            \vdots\\
            \expectation[X_n]
        \end{pmatrix}.
    \]
\end{snippetdefinition}

\begin{snippetdefinition}{covariance-matrix-definition}{Covariance matrix}
    Let \((\Omega,\mathcal F,\mathbb P)\) be a \probabilityspace, and let
    \(X=(X_1,\ldots,X_n)\) be a \randomvector in \(\realnumbers^n\). If each
    component has finite second moment, the \emph{covariance matrix} of \(X\)
    is
    \[
        Q_X=(q_{ij})_{i,j=1}^n,
        \qquad
        q_{ij}=\covariance(X_i,X_j).
    \]
\end{snippetdefinition}

\begin{snippetproposition}{covariance-matrix-basic-properties}{Basic properties of the covariance matrix}
    Let \((\Omega,\mathcal F,\mathbb P)\) be a \probabilityspace, and let
    \(X=(X_1,\ldots,X_n)\) be a \randomvector in \(\realnumbers^n\) whose
    components have finite second moments. Then:
    \begin{enumerate}
        \item \(Q_X\) is symmetric;
        \item
            \[
                \operatorname{diag}(Q_X)
                =
                \bigl(\variance(X_1),\ldots,\variance(X_n)\bigr);
            \]
        \item for every \(u\in\realnumbers^n\),
            \[
                u^TQ_Xu=\variance(u^TX)\geq0.
            \]
    \end{enumerate}
\end{snippetproposition}

\begin{snippetproof}{covariance-matrix-basic-properties-proof}{covariance-matrix-basic-properties}{Basic properties of the covariance matrix}
    Symmetry follows from
    \(\covariance(X_i,X_j)=\covariance(X_j,X_i)\). The diagonal identity
    follows from \(\covariance(X_i,X_i)=\variance(X_i)\). Finally, using
    bilinearity of covariance,
    \[
        u^TQ_Xu
        =
        \sum_{i=1}^n\sum_{j=1}^n u_i u_j\covariance(X_i,X_j)
        =
        \variance\left(\sum_{i=1}^n u_iX_i\right)
        =
        \variance(u^TX)\geq0.
    \]
\end{snippetproof}

\begin{snippetproposition}{affine-transformation-mean-covariance}{Mean vector and covariance matrix under affine transformations}
    Let \((\Omega,\mathcal F,\mathbb P)\) be a \probabilityspace, and let
    \(X\) be a \randomvector in \(\realnumbers^n\) whose components have
    finite second moments. Let \(A\in\realnumbers^{m\times n}\), let
    \(b\in\realnumbers^m\), and define
    \[
        Y=AX+b.
    \]
    Then
    \[
        \expectation[Y]=A\expectation[X]+b,
        \qquad
        Q_Y=AQ_XA^T.
    \]
\end{snippetproposition}

\begin{snippetproof}{affine-transformation-mean-covariance-proof}{affine-transformation-mean-covariance}{Mean vector and covariance matrix under affine transformations}
    Linearity of the \expectedvalue gives
    \[
        \expectation[Y]=A\expectation[X]+b.
    \]
    Since translating by \(b\) does not change covariance,
    \[
        \covariance(Y_r,Y_s)
        =
        \covariance\left(\sum_{i=1}^n a_{ri}X_i,
        \sum_{j=1}^n a_{sj}X_j\right)
        =
        \sum_{i=1}^n\sum_{j=1}^n
        a_{ri}a_{sj}\covariance(X_i,X_j),
    \]
    which is precisely the \((r,s)\)-entry of \(AQ_XA^T\).
\end{snippetproof}

\section{Gaussian vectors}

\begin{snippetdefinition}{standard-gaussian-vector-definition}{Standard Gaussian vector}
    Let \((\Omega,\mathcal F,\mathbb P)\) be a \probabilityspace, and let
    \(X=(X_1,\ldots,X_n)\) be a \randomvector in \(\realnumbers^n\). The
    vector \(X\) is a \emph{standard Gaussian vector}, written
    \[
        X\sim\mathcal N(0,I_n),
    \]
    if \(X\) is an \absolutelycontinuousrandomvector with density
    \[
        f_X(x_1,\ldots,x_n)
        =
        \frac{1}{(2\pi)^{n/2}}
        e^{-\frac12(x_1^2+\cdots+x_n^2)}.
    \]
\end{snippetdefinition}

\begin{snippetproposition}{standard-gaussian-vector-properties}{Properties of the standard Gaussian vector}
    Let \(X=(X_1,\ldots,X_n)\) be a \standardgaussianvector. Then
    \(X_1,\ldots,X_n\) are \independentrandomvariables,
    \[
        X_i\sim\mathcal N(0,1)
        \qquad
        \forall i\in\{1,\ldots,n\},
    \]
    and
    \[
        \expectation[X]=0,
        \qquad
        Q_X=I_n.
    \]
\end{snippetproposition}

\begin{snippetproof}{standard-gaussian-vector-properties-proof}{standard-gaussian-vector-properties}{Properties of the standard Gaussian vector}
    The density factorizes as
    \[
        f_X(x_1,\ldots,x_n)
        =
        \prod_{i=1}^n
        \frac{1}{\sqrt{2\pi}}e^{-\frac{x_i^2}{2}}.
    \]
    Hence the density factorization criterion gives independence, and each
    marginal is \(\mathcal N(0,1)\). The one-dimensional Gaussian formulas
    give \(\expectation[X_i]=0\) and \(\variance(X_i)=1\). Since independent
    variables with finite second moments have covariance \(0\), the
    covariance matrix is \(I_n\).
\end{snippetproof}

\begin{snippetdefinition}{gaussian-vector-definition}{Gaussian vector}
    Let \((\Omega,\mathcal F,\mathbb P)\) be a \probabilityspace, let
    \(m\in\realnumbers^n\), and let \(Q\in\realnumbers^{n\times n}\) be a
    \symmetricmatrix and \positivedefinitematrix, meaning that
    \[
        u^TQu>0
        \qquad
        \forall u\in\realnumbers^n\difference\{0\}.
    \]
    A \randomvector \(X\colon\Omega\fromto\realnumbers^n\) is a
    \emph{Gaussian vector} with parameters \(m\) and \(Q\), written
    \[
        X\sim\mathcal N(m,Q),
    \]
    if \(X\) is an \absolutelycontinuousrandomvector with density
    \[
        f_X(x)
        =
        \frac{1}{(2\pi)^{n/2}\sqrt{\det(Q)}}
        e^{-\frac12 (x-m)^TQ^{-1}(x-m)}.
    \]
\end{snippetdefinition}

\begin{snippetlemma}{positive-definite-square-root-for-gaussians}{Square root of a positive definite matrix}
    Let \(Q\in\realnumbers^{n\times n}\) be a \symmetricmatrix and
    \positivedefinitematrix. There exists a unique \symmetricmatrix and
    \positivedefinitematrix \(Q^{1/2}\) such that
    \[
        Q^{1/2}Q^{1/2}=Q.
    \]
    Moreover, \(Q^{1/2}\) is invertible and
    \[
        \det(Q^{1/2})=\sqrt{\det(Q)}.
    \]
\end{snippetlemma}

\begin{snippetproof}{positive-definite-square-root-for-gaussians-proof}{positive-definite-square-root-for-gaussians}{Square root of a positive definite matrix}
    By the real \spectraltheorem, there are an \orthogonalmatrix \(P\) and
    positive numbers \(\lambda_1,\ldots,\lambda_n\) such that
    \[
        Q=P\operatorname{diag}(\lambda_1,\ldots,\lambda_n)P^T.
    \]
    Define
    \[
        Q^{1/2}
        =
        P\operatorname{diag}(\sqrt{\lambda_1},\ldots,\sqrt{\lambda_n})P^T.
    \]
    This matrix is symmetric, positive definite, and squares to \(Q\). The
    determinant formula follows by multiplying the eigenvalues. Uniqueness
    follows because a symmetric positive definite square root must have the
    same eigenspaces and the positive square roots of the eigenvalues.
\end{snippetproof}

\begin{snippettheorem}{gaussian-vector-standard-characterization}{Characterization through the standard Gaussian vector}
    Let \(m\in\realnumbers^n\), and let \(Q\in\realnumbers^{n\times n}\) be a
    \symmetricmatrix and \positivedefinitematrix. Let \(X\) be a
    \randomvector in \(\realnumbers^n\). Then
    \[
        X\sim\mathcal N(m,Q)
    \]
    if and only if
    \[
        X\eqdist Q^{1/2}Z+m,
        \qquad
        Z\sim\mathcal N(0,I_n).
    \]
\end{snippettheorem}

\begin{snippetproof}{gaussian-vector-standard-characterization-proof}{gaussian-vector-standard-characterization}{Characterization through the standard Gaussian vector}
    Suppose first that \(Y=Q^{1/2}Z+m\), where \(Z\sim\mathcal N(0,I_n)\).
    Since \(Q^{1/2}\) is invertible, the \affinedensitytransformation gives
    \[
        f_Y(y)
        =
        \frac{1}{\det(Q^{1/2})}
        \frac{1}{(2\pi)^{n/2}}
        e^{-\frac12\left\|Q^{-1/2}(y-m)\right\|^2}.
    \]
    Since \(Q^{-1/2}\) is symmetric,
    \[
        \left\|Q^{-1/2}(y-m)\right\|^2
        =
        (y-m)^TQ^{-1}(y-m),
    \]
    and \(\det(Q^{1/2})=\sqrt{\det(Q)}\). Thus \(Y\sim\mathcal N(m,Q)\).
    Therefore any \(X\eqdist Y\) also satisfies \(X\sim\mathcal N(m,Q)\).

    Conversely, suppose \(X\sim\mathcal N(m,Q)\). Let
    \[
        Z=Q^{-1/2}(X-m).
    \]
    Applying the same affine density transformation gives
    \[
        f_Z(z)=\frac{1}{(2\pi)^{n/2}}e^{-\frac12 z^Tz}.
    \]
    Hence \(Z\sim\mathcal N(0,I_n)\), and
    \(X=Q^{1/2}Z+m\).
\end{snippetproof}

\begin{snippettheorem}{gaussian-vector-mean-covariance}{Parameters of a Gaussian vector}
    Let \(m\in\realnumbers^n\), let \(Q\in\realnumbers^{n\times n}\) be a
    \symmetricmatrix and \positivedefinitematrix, and let
    \(X\sim\mathcal N(m,Q)\). Then
    \[
        \expectation[X]=m,
        \qquad
        Q_X=Q.
    \]
\end{snippettheorem}

\begin{snippetproof}{gaussian-vector-mean-covariance-proof}{gaussian-vector-mean-covariance}{Parameters of a Gaussian vector}
    By the characterization through the \standardgaussianvector,
    \[
        X\eqdist Q^{1/2}Z+m,
        \qquad
        Z\sim\mathcal N(0,I_n).
    \]
    Equality in law preserves the \expectedvector and \covariancematrix. Thus,
    using the affine transformation formulas,
    \[
        \expectation[X]
        =
        Q^{1/2}\expectation[Z]+m
        =
        m,
    \]
    and
    \[
        Q_X
        =
        Q^{1/2}I_n(Q^{1/2})^T
        =
        Q.
    \]
\end{snippetproof}

\section{Stability}

\begin{snippettheorem}{invertible-affine-transformation-gaussian-vector}{Invertible affine transformations of Gaussian vectors}
    Let \(X\sim\mathcal N(m,Q)\) be a \gaussianvector in
    \(\realnumbers^n\). Let \(A\in\realnumbers^{n\times n}\) be invertible,
    let \(b\in\realnumbers^n\), and define
    \[
        Y=AX+b.
    \]
    Then
    \[
        Y\sim\mathcal N(Am+b,AQA^T).
    \]
\end{snippettheorem}

\begin{snippetproof}{invertible-affine-transformation-gaussian-vector-proof}{invertible-affine-transformation-gaussian-vector}{Invertible affine transformations of Gaussian vectors}
    By the \affinedensitytransformation,
    \[
        f_Y(y)
        =
        \frac{1}{|\det A|}
        f_X(A^{-1}(y-b)).
    \]
    Substituting the Gaussian density of \(X\) and using
    \[
        \det(AQA^T)=(\det A)^2\det(Q),
        \qquad
        (AQA^T)^{-1}=(A^T)^{-1}Q^{-1}A^{-1},
    \]
    gives
    \[
        f_Y(y)
        =
        \frac{1}{(2\pi)^{n/2}\sqrt{\det(AQA^T)}}
        e^{-\frac12 (y-Am-b)^T(AQA^T)^{-1}(y-Am-b)}.
    \]
    Hence \(Y\sim\mathcal N(Am+b,AQA^T)\).
\end{snippetproof}

\begin{snippetproposition}{linear-combination-standard-gaussian}{Linear combinations of independent standard Gaussians}
    Let \(Z\sim\mathcal N(0,I_n)\), and let
    \(a\in\realnumbers^n\difference\{0\}\). Then
    \[
        a^TZ\sim\mathcal N(0,a^Ta).
    \]
\end{snippetproposition}

\begin{snippetproof}{linear-combination-standard-gaussian-proof}{linear-combination-standard-gaussian}{Linear combinations of independent standard Gaussians}
    Choose an \orthogonalmatrix \(O\) whose first row is \(a^T/\sqrt{a^Ta}\).
    Since orthogonal transformations preserve the Euclidean norm and have
    determinant of absolute value \(1\), the vector \(OZ\) is again standard
    Gaussian. Therefore the first component of \(OZ\) is \(\mathcal N(0,1)\).
    But
    \[
        (OZ)_1=\frac{a^TZ}{\sqrt{a^Ta}},
    \]
    so \(a^TZ\sim\mathcal N(0,a^Ta)\).
\end{snippetproof}

\begin{snippettheorem}{gaussian-components-independence-diagonal-covariance}{Independence of Gaussian components}
    Let
    \[
        X=(X_1,\ldots,X_n)\sim\mathcal N(m,Q),
    \]
    where \(m=(m_1,\ldots,m_n)^T\). Then \(X_1,\ldots,X_n\) are
    \independentrandomvariables if and only if \(Q\) is diagonal. More
    explicitly, if
    \[
        Q=\operatorname{diag}(\lambda_1,\ldots,\lambda_n),
        \qquad
        \lambda_i>0,
    \]
    then
    \[
        X_i\sim\mathcal N(m_i,\lambda_i)
        \qquad
        \forall i\in\{1,\ldots,n\},
    \]
    and the components are independent.
\end{snippettheorem}

\begin{snippetproof}{gaussian-components-independence-diagonal-covariance-proof}{gaussian-components-independence-diagonal-covariance}{Independence of Gaussian components}
    If \(Q=\operatorname{diag}(\lambda_1,\ldots,\lambda_n)\), then the
    density of \(X\) is
    \[
        f_X(x_1,\ldots,x_n)
        =
        \prod_{i=1}^n
        \frac{1}{\sqrt{2\pi\lambda_i}}
        e^{-\frac{(x_i-m_i)^2}{2\lambda_i}}.
    \]
    Hence the joint density factorizes into one-dimensional Gaussian
    densities, so the components are independent and
    \(X_i\sim\mathcal N(m_i,\lambda_i)\).

    Conversely, if the components are independent, then
    \(\covariance(X_i,X_j)=0\) for \(i\neq j\). Since the covariance matrix of
    \(X\) is \(Q\), all off-diagonal entries of \(Q\) are zero.
\end{snippetproof}

\begin{snippetcorollary}{sum-of-two-jointly-gaussian-variables}{Sum of two jointly Gaussian variables}
    Let
    \[
        (X,Y)\sim\mathcal N(m,Q),
        \qquad
        m=
        \begin{pmatrix}
            m_X\\
            m_Y
        \end{pmatrix},
        \qquad
        Q=
        \begin{pmatrix}
            q_{11} & q_{12}\\
            q_{12} & q_{22}
        \end{pmatrix}.
    \]
    Then
    \[
        X+Y\sim\mathcal N(m_X+m_Y,q_{11}+q_{22}+2q_{12}).
    \]
    Equivalently,
    \[
        X+Y
        \sim
        \mathcal N\left(
            \expectation[X]+\expectation[Y],
            \variance(X)+\variance(Y)+2\covariance(X,Y)
        \right).
    \]
\end{snippetcorollary}

\begin{snippetproof}{sum-of-two-jointly-gaussian-variables-proof}{sum-of-two-jointly-gaussian-variables}{Sum of two jointly Gaussian variables}
    Let \(\ell=(1,1)^T\). By the standard characterization,
    \[
        \begin{pmatrix}
            X\\
            Y
        \end{pmatrix}
        \eqdist
        m+Q^{1/2}Z,
        \qquad
        Z\sim\mathcal N(0,I_2).
    \]
    Therefore
    \[
        X+Y
        \eqdist
        \ell^Tm+\ell^TQ^{1/2}Z.
    \]
    By the linear-combination result, the centered term is Gaussian with
    variance
    \[
        \ell^TQ^{1/2}(Q^{1/2})^T\ell
        =
        \ell^TQ\ell
        =
        q_{11}+q_{22}+2q_{12}.
    \]
    The mean is \(\ell^Tm=m_X+m_Y\).
\end{snippetproof}

\section{Degenerate Gaussian vectors}

\begin{snippetnote}{classical-degenerate-gaussian-vector-note}{Degenerate Gaussian vectors}
    There is a broader convention in which a \randomvector
    \(X\colon\Omega\fromto\realnumbers^n\) is called Gaussian if, for every
    \(u\in\realnumbers^n\), the real-valued random variable \(u^TX\) has
    either a \gaussiandistribution or a \diracmeasure as its law. In that
    convention the \covariancematrix is only positive semidefinite, and the
    law need not be
    absolutely continuous with respect to \(n\)-dimensional
    \lebesguemeasure.
\end{snippetnote}

\begin{snippetexample}{degenerate-gaussian-vector-example}{Degenerate Gaussian vector}
    Let \(Z\sim\mathcal N(0,1)\), and define
    \[
        X=(Z,0).
    \]
    Then \(X\) is Gaussian in the degenerate sense. Its \covariancematrix is
    \[
        Q_X=
        \begin{pmatrix}
            1 & 0\\
            0 & 0
        \end{pmatrix}.
    \]
    The law of \(X\) is concentrated on the line
    \[
        \{(x,y)\in\realnumbers^2\suchthat y=0\},
    \]
    which has two-dimensional \lebesguemeasure zero. Hence \(X\) is not an
    \absolutelycontinuousrandomvector in \(\realnumbers^2\).
\end{snippetexample}

\end{document}
